\section{Interpolation with two arguments}\label{sec:interpolation}

This section proves a polynomial height bound from common analytic
relations at a weighted collection of places.  The relations involve a
small algebraic parameter $a$, auxiliary algebraic coordinates $U$, and
solutions of fixed ordinary differential systems.  A second argument may
be $aR$, where $R$ is a coordinate of $U$.  The conclusion will bound a
height parameter $h$ polynomially in a complexity parameter $d$, provided
the relations are not all identities on any of the algebraic branches
used in a dimension descent.

\subsection{Weighted absolute heights}

For a tuple $z=(z_i)$ over a number field $F$, write
$\norm z_v=\max_i|z_i|_v$ and put
\[
 \ha{z}=\sum_v w_v\log\max(1,\norm{z}_v),
 \qquad w_v=\frac{[F_v:\Q_v]}{[F:\Q]}.
\]
At the infinite places the absolute value is the usual, unsquared one;
thus a complex place has weight $2/[F:\Q]$.  The same notation for a
polynomial, or a finite list of polynomials, means the height of its
complete coefficient tuple \emph{with the additional coordinate $1$}.
It therefore bounds the scale of the coefficients.  For each fixed
number of variables,
\begin{equation}\label{eq:polynomial-evaluation}
 \log\abs{G(z)}_v
 \leq \log\max(1,\norm{\operatorname{coeff}G}_v)
       +\deg(G)\log\max(1,\norm z_v)
       +\eps_v\log\#\operatorname{mon}(G),
\end{equation}
where $\eps_v$ is $1$ at an infinite place and $0$ otherwise.
In particular, the sum of the last terms is controlled without counting
the finite places.

Occasionally the selected local data distinguish embeddings inducing the
same place.  We then use embeddings $F\hookrightarrow\C$ or
$F\hookrightarrow\C_p$, each with weight $1/[F:\Q]$, which gives the
same height formula.  A selection of embeddings is allowed.  On replacing
$F$ by a finite extension, we always take \emph{all} extensions of each
selected embedding.  The sum of their weights is its original weight.
All selected sums below are consequently unchanged by such an extension.
No bound for the degree of a normal closure or of an auxiliary coefficient
field will be used.

\subsection{The interpolation statement}

In the statement, the finite list of curves, connections, ordinary
centres, rational frames regular at the centres, and algebraic chart
germs is fixed.  Its ranks and all polynomial-combination degrees are
fixed.  An allowed analytic expression is a polynomial combination of
the algebraic coordinates $(a,U)$, these chart germs, and entries of the
horizontal solution matrices, evaluated at $a$ and possibly at $aR$,
where $R$ is one of the coordinates of $U$.  Constant summands are
allowed.  The total degree in all polynomial factors, including $a$
and every coordinate of $U$, is bounded by a fixed constant.  Only
substitution into the fixed germs produces Taylor coefficients of
growing degree.  The scalar coefficients of the combination may be
algebraic and vary with the target.

\begin{theorem}[Common-relation interpolation]\label{thm:interpolation}
Fix the analytic data just described, a fixed affine ambient space with
coordinates $(a,U)$, and an algebraic ambient locus $X$ in that space.
Let $h,d\geq2$, let $\Psi(h)=1+(\log h)^c$ for a fixed $c\geq1$, and
let $z_*=(a_*,U_*)\in X(\Qbar)$ with $a_*\ne0$.  Suppose that the
following bounds hold, with fixed polynomials $Q_i$ having nonnegative
coefficients.
\begin{enumerate}[label=\textup{(\roman*)}]
\item The ambient locus has equations of degree at most $Q_0(d)$ and
affine height at most $Q_0(d)\Psi(h)$, and
$\ha{z_*}\leq Q_0(d)h$.
\item There is a finite weighted selection $\mathcal V$ of places or
embeddings for which $\abs{a_*}_v<1$ and
\begin{equation}\label{eq:interpolation-smallness}
 S_{\rm small}:=\sum_{v\in\mathcal V}w_v
               \log\frac1{\abs{a_*}_v}
 \geq\frac{h}{Q_1(d)},\qquad
 \sum_{v\in\mathcal V}w_v\log\max(1,\norm{U_*}_v)
 \leq Q_2(d)\Psi(h).
\end{equation}
If the second argument occurs, $\abs{R_*}_v=1$ for these local data.
\item A common list of fixed length $s$ of allowed expressions
\[
 E_\nu(a,U)=\sum_{j\geq0}a^j e_{\nu j}(U),\qquad 1\leq\nu\leq s,
\]
vanishes at $z_*$ at every selected local datum.  Its coefficients satisfy
\[
 \deg e_{\nu j}\leq C(1+j),\qquad
 \ha{(e_{\nu j})_{\nu,\,0\leq j\leq l}}
 \leq Q_3(l,d)\Psi(h).
\]
There are $b_v\geq0$ with
$\sum_{v\in\mathcal V}w_v b_v\leq Q_4(d)\Psi(h)$ such that the
convergent tails at the target satisfy, for every $l\geq0$,
\begin{equation}\label{eq:interpolation-tail}
 \left|\sum_{j\geq l}a_*^j e_{\nu j}(U_*)\right|_v
 \leq \abs{a_*}_v^l\exp((1+l)b_v).
\end{equation}
\item If $W\subseteq X$ is an irreducible algebraic variety containing
$z_*$ on which $a$ is nonconstant, then on every normalization branch
above the generic point of every component of $W\cap V(a)$ the
expressions $E_1,\ldots,E_s$ are not all identically zero.  The
expressions use their specified ordinary chart germs.
\end{enumerate}
Then $h\leq C_1d^{C_2}$, with constants depending only on the fixed data
and the polynomials in the hypotheses.
\end{theorem}

The hypothesis in (iv) includes branches on proper subvarieties of $X$.
The proof only uses the varieties encountered in the descent below, so
verification on those varieties suffices.  In the application this is
the role of \Cref{prop:nonidentity}.

Here is the descent mechanism.  Suppose that $h$ is larger than a
polynomial in $d$, and let $W\subseteq X$ be an irreducible variety
containing the target.  When $a$ is nonconstant and $W$ meets $a=0$,
we pass to a graph $W'$ that is isomorphic to $W$ over $a\ne0$.
On this graph we construct a polynomial $P$ with
\[
 P\notin I(W'),\qquad P(z_*')=0,
\]
where $z_*'$ is the lifted target.  The intersection with $P=0$ then
projects to a smaller-dimensional variety containing $z_*$.  We discard
the graph coordinate after each step, so the ambient dimension stays
fixed.

The two requirements on $P$ have different sources.  Nonidentity in
(iv), together with an ordinary-point zero estimate, gives a graph whose
special fibre satisfies a codimension-two condition after one further
equation is imposed.  Counting coefficients layer by layer in powers of
$a$ then gives a polynomial nonzero on $W'$ whose target value is bounded
locally by $|a_*|_v^n$ times controlled factors.  The ratio of $n$ to its
degree $L$ can exceed any prescribed polynomial in $d$, while both $n$
and $L$ remain polynomially bounded.  At the selected places, the factor
$|a_*|_v^n$ contributes $-nS_{\rm small}$ to the logarithm of its value.
At the other places, its degree costs at most $L$ times the global height
of the lifted target.  For a sufficiently large ratio
$n/L$, the product formula forces $P(z_*')=0$.  The remaining errors
are polynomial multiples of $\Psi(h)$ and are absorbed when $h$ is large.

\subsection{Effective algebra in fixed dimension}

The descent needs two kinds of algebraic bounds: equations for each
smaller variety must retain controlled degree and height, and the
interpolation polynomial must have a membership identity with controlled
coefficients.  We record both bounds here.  Throughout this subsection,
``controlled'' means polynomial in the input degrees and linear in
$1+$ the joint logarithmic coefficient height, with ambient dimension
fixed.

In a fixed number of variables, ideals with generators of degree at most $D$ have
Gr\"obner bases of degree polynomial in $D$, for any prescribed monomial
order \cite[Corollary~8.3]{Dube1990}.  We may first discard generators
linearly dependent on the others, leaving polynomially many.  If a
polynomial belongs to such an
ideal, the degrees in a membership representation are polynomial in
$D$ and the degree of that polynomial
\cite[Theorem~3.4]{Aschenbrenner2004}.  This is a bound for arbitrary
ideal membership.

These degree bounds also give coefficient bounds from the original
input coefficients.  Let $\delta$ be the Gr\"obner-basis degree bound,
and choose the minimal monomial generators $M_\alpha$ of the initial
ideal.  Each is the leading monomial of a polynomial of degree at most
$\delta$ in the ideal.  Seek such a polynomial in the form
\[
 H_\alpha=M_\alpha+
 \sum_{\substack{N<M_\alpha\\\deg N\leq\delta}}c_{\alpha N}N,
 \qquad H_\alpha=\sum_i A_{\alpha i}G_i.
\]
The membership bound supplies a polynomial bound for the degrees of
the $A_{\alpha i}$.  The displayed coefficient equations are therefore
a solvable affine linear system of polynomial size whose entries are
only the original coefficients and $0,\pm1$.  A solution obtained by
nonzero minors bounds the coefficients of $H_\alpha$ and its
representation simultaneously.  The resulting $H_\alpha$ form a
Gr\"obner basis, since their leading monomials generate the initial
ideal.  Thus no bound for coefficients of a previously chosen basis
is presumed.

If the joint affine height of the input is $H$, a determinant of size
$s$ has height at most $sH+s\log s$.  Normalizing a nonzero coordinate
and concatenating polynomially many coefficient lists incurs only a
polynomial factor.  Consequently all these affine coefficient heights
are at most $Q(D)(H+1)$: the dependence on the logarithmic input height
is linear.  The same argument treats elimination and saturation; for
example saturation by $a$ is computed by eliminating $z$ after
adjoining $za-1$.

For varieties, the degree and arithmetic B\'ezout inequalities give the
following consequence: intersections, closures of coordinate projections,
and choices of irreducible components have degrees polynomial in the
input degrees and Chow heights at most a polynomial in the input degrees times $1+$
the joint input height.  Here and below the ambient dimensions are fixed.
We use the arithmetic B\'ezout inequality
\cite[Corollary~2.11]{KPS2001}, the projection bound
\cite[Lemma~2.6]{KPS2001}, and additivity and nonnegativity of cycle
heights \cite[Section~1.2.4]{KPS2001}.  The comparison of Chow
coefficient height with this cycle height is
\cite[Lemma~1.1, inequality~(1.2), and Section~1.2.4]{KPS2001}.
One may work throughout with Chow forms, normalizing one nonzero
coefficient to $1$.  The affine height of the normalized coefficient
tuple then equals its projective height, by the product formula.

It is enough to recover equations \emph{set-theoretically} from a Chow
form.  Indeed, if $Z\subset\mathbb P^N$ has dimension $r$ and Chow form
$\mathcal C_Z$, substitute the $r+1$ hyperplanes
\[
 \sum_{j=1}^N T_{ij}(X_j-x_jX_0)=0\qquad(0\leq i\leq r)
\]
in $\mathcal C_Z$ and take its coefficients as a polynomial in the
$T_{ij}$.  These coefficients vanish at $x$ precisely when
$[1:x]\in Z$: outside $Z$, the hyperplanes through that point form a
base-point-free system on $Z$, and $r+1$ general members have empty
intersection on $Z$.  The resulting equations have degree at most
$(r+1)\deg Z$ and controlled affine coefficient heights by expansion.
Their number is polynomial in $\deg Z$.

Thus every variety used below admits an ideal $\mathfrak b$ of controlled
generators with
\begin{equation}\label{eq:controlled-support}
 V(\mathfrak b)=W,\qquad \mathfrak b\subseteq I(W).
\end{equation}
We need not compute generators of the radical $I(W)$.  The distinction
will be useful: ideal membership in $\mathfrak b$ suffices for evaluation
on $W$, while nonvanishing on $W$ will be imposed separately.

\subsection{An ordinary-point zero estimate}

The differential method compares successive derivatives in a
finite-dimensional row space; see the classical multiplicity estimate
of Bertrand and Beukers \cite[Theorem~1]{BertrandBeukers1985}.
We give the argument here because the algebraic variety on which the
system is restricted varies in the descent, and the bound must be
polynomial in its degree while allowing the second argument $aR$.

\begin{lemma}[Uniform ordinary-point order bound]\label{lem:zero-estimate}
For the fixed analytic data of \Cref{thm:interpolation} there is a
polynomial $M(D)$ with the following property.  Let $W$ be an irreducible
affine variety of degree at most $D$ on which $a$ is nonconstant.  On a
branch above the generic point of a component of $V_W(a)$, every allowed
expression $E$ which is not identically zero satisfies
\begin{equation}\label{eq:divisorial-zero-estimate}
 \operatorname{ord}_{\mathfrak p}E
 <M(D)\operatorname{ord}_{\mathfrak p}a.
\end{equation}
The bound is independent of the heights of the coefficients of $W$ and
of the expression.
\end{lemma}

\begin{proof}
We first preserve the branch order on a curve of controlled degree, then
use the ordinary differential equation to bound that order independently
of coefficient heights.

\emph{Reduction to a curve.}  Write $k=\dim W$, and
let $\widetilde D$ be the normalization divisor defining the branch.
At a general point $p$ of $\widetilde D$, both the normalization and
$\widetilde D$ are smooth, and the finite map from $\widetilde D$ to
its image $D_0$ is \'etale at $p$ after restricting to dense open
subsets, by separability in characteristic zero.  Choose $k-1$ general affine linear functions
whose differentials restrict to a basis of the cotangent space of
$D_0$ at the image of $p$.  Their pullbacks restrict to a basis on
$\widetilde D$.  Their common level set through $p$ is consequently
a smooth curve germ transverse to $\widetilde D$ in the normalization.
The downstairs linear section is proper and has total projective
degree at most $D$.  The transverse germ leaves the nonnormal locus,
so its map to its image curve is generically the normalization
isomorphism.  It is thus the corresponding smooth germ on the
normalization of an irreducible component of that linear section.

This slice preserves formal orders as well.  If
$\operatorname{ord}_{\widetilde D}a=e>0$ and
$\operatorname{ord}_{\widetilde D}E=N$, only the terms with $j\leq N/e$
contribute to the leading coefficient.  After deleting a proper closed
subset of $\widetilde D$, write $a=t^ev$ and that finite truncation as
$t^Nu$, where $t$ is a local equation of $\widetilde D$ and $u,v$ are
units.  On the transverse curve $t$ is a uniformizer.  Hence the
orders are exactly $e$ and $N$ after restriction.  All these conditions
hold at a general slice point.

Let $T_0$ be the smooth projective normalization of this sliced curve.
Its genus is polynomial in $D$, by a general birational plane
projection and the arithmetic genus of a plane curve.  Every ambient
coordinate on $T_0$ is a quotient $X_i/X_0$, so its polar divisor has
degree at most $D$.  In particular the maps defined by $a$ and $aR$
have degrees at most $D$ and $2D$, respectively.

Adjoin the points on the fixed curves specified by the algebraic chart
germs.  Constant argument maps contribute constant chart factors and
require no cover; in what follows index only the nonconstant argument
maps.  If the corresponding fixed maps to $\mathbb P^1$ have degrees $d_i$ and
branch sets $B_i$, including infinity when it is a branch value, the
chosen normalized fibre-product component $T\to T_0$ has degree at
most $N_0=\prod_i d_i$.  It is \'etale outside
$\bigcup_i\phi_i^{-1}(B_i)$, where $\phi_i$ is $a$ or $aR$.
This set has at most $\sum_i\#B_i\deg\phi_i=O(D)$ points, and the
ramification contribution above each is at most $N_0$.  Thus
Riemann--Hurwitz gives
\[
 2g(T)-2\leq N_0\max\{0,2g(T_0)-2\}+O(D).
\]
The degrees of the maps
from $T$ to the fixed curves are bounded by $N_0\deg\phi_i$, so all
pulled-back fixed divisors have polynomial degree.  Poles of $aR$ are
included in this calculation through the fibre above infinity.  All
ramification indices are bounded as well; any further ramification
multiplies the two orders in
\eqref{eq:divisorial-zero-estimate} by the same factor.

\emph{The ordinary differential system.}  Let $p\in T$ be the point of
the specified branch.  In direct sums and tensor products of the fixed
pulled-back systems the expression takes the form
\[
 E=r_0y,
\]
where $y$ is a horizontal column of a fixed rank $N$ and $r_0$ is a
rational row.  Polynomial algebraic factors are absorbed in $r_0$.
A constant map to one of the fixed curves contributes a constant
horizontal factor.  The rational frames here are the pullbacks of the
fixed frames; they are regular and invertible at $p$.  In fact $a=0$
and $R$ is regular at the chosen affine divisor, so both arguments
$a$ and $aR$ are at their ordinary centres.  The connection and $r_0$
are therefore regular at $p$.  Their global polar divisors have degrees
polynomial in $D$, by the bounds on the degrees of the pullback maps
and of the coordinate functions.

Choose a rational function $z$ on $T$ which is a uniformizer at $p$ and
whose polar divisor has degree at most $2g(T)+1$.  Such a function follows
from Riemann--Roch by prescribing its first two jets at $p$ and allowing
a pole at another point.  For $g(T)=0$ a degree-one function suffices.
Put $\delta=d/dz$, and write the horizontal equation as $\delta y=Ay$.
The poles of the rational vector field $\delta$ are bounded in degree by
the zero divisor of $\dd z$, hence by $2g(T)-2+2\deg z$.  Differentiating
a rational function adds at most its reduced polar divisor and the
polar divisor of this vector field.  It follows that $A$, $r_0$, and
any fixed number of their derivatives have polar degrees polynomial in
$D$.  They remain regular at $p$; using a uniformizer rather than
dividing by $\dd a$ accommodates arbitrary ramification of $a$.

Set
\[
 r_{i+1}=\delta r_i+r_iA.
\]
Then $\delta^iE=r_iy$.  The successive rows become dependent after at
most $N$ steps.  If $r_\rho$ is the first row dependent on its
predecessors, their span
$M=\langle r_0,\ldots,r_{\rho-1}\rangle_{\C(T)}$ is stable under the
dual connection: differentiating the dependence shows that every later
row belongs to it.  Here $\rho\leq N$.

Let $\mathcal O=\mathcal O_{T,p}$ and saturate the common kernel of
these rows in $\mathcal O^N$.  Saturation makes it a subbundle locally.
It is stable under the ordinary connection, since its derivative is
regular and belongs to the kernel generically.  The quotient has an
ordinary connection as well.  Explicitly, if $v$ is a local section
annihilated by every row in $M$, stability of $M$ under the dual
connection gives $r(\nabla_\delta v)=0$ for every $r\in M$.
The ordinary connection makes $\nabla_\delta v$ regular, so it belongs
to the saturated kernel.  The kernel and its quotient are free over
the local discrete valuation ring, and the induced quotient connection
has no pole.  The image $\bar y$ of $y$ in this
quotient is a nonzero horizontal section: a zero image would give
$r_0y=E=0$.  Conversely, if $E$ were identically zero, all its
derivatives $r_i y$ would be zero and $y$ would lie in the kernel.
The nonzero horizontal quotient section has
nonzero reduction at $p$: the coefficient recurrence of an ordinary
linear differential equation forces a horizontal section with zero
initial value to be zero.

The integral dual of this free quotient is exactly the saturated row
lattice $\mathcal L=M\cap\mathcal O^N$: its rows are the integral
functionals annihilating the common kernel.  Since $\bar y$ has nonzero
reduction, one such functional $\lambda\in\mathcal L$ has
$\lambda y$ a unit.  Choose a nonzero $\rho\times\rho$ minor $\Delta$
of the matrix with rows $r_0,\ldots,r_{\rho-1}$.  All these rows are
regular at $p$.  Cramer's rule on the selected columns gives
\[
 \Delta\lambda=\sum_{i=0}^{\rho-1}c_ir_i,
 \qquad c_i\in\mathcal O.
\]
Hence at least one $i<\rho$ satisfies
$\operatorname{ord}_p(r_iy)\leq\operatorname{ord}_p\Delta$.
This is the elementary-divisor bound for the row lattice, expressed
without choosing a new global frame.  A nonzero rational function on
$T$ has equal total zero and pole degrees, so
$\operatorname{ord}_p\Delta$ is bounded by its polar degree, which is
polynomial in $D$.  Finally, since $z$ is a uniformizer,
\[
 \operatorname{ord}_p E
 \leq\operatorname{ord}_p(\delta^iE)+i
 \leq\operatorname{ord}_p\Delta+N-1.
\]
Increasing the resulting polynomial by $1$ proves
\eqref{eq:divisorial-zero-estimate}, since
$\operatorname{ord}_p a\geq1$.  This construction used only degrees,
genera, ramification and fixed ranks; it used no arithmetic height.
\end{proof}

\subsection{The graph and its special fibre}

\begin{lemma}[Dimension of the constrained graph fibre]\label{lem:graph-fiber}
Let $W$ be an irreducible affine variety of dimension $k\geq1$ with
coordinate $a$ nonconstant, let $f$ be regular on $W$, and let $m\ge1$
be an integer.  Suppose
\[
 \operatorname{ord}_{\mathfrak p}f
 <m\operatorname{ord}_{\mathfrak p}a
\]
on every normalization branch over the generic components of $V_W(a)$.
Let $W'$ be the affine graph closure of $q=f/a^m$ over $W\setminus V(a)$.
For every regular function $g$ on $W$,
\begin{equation}\label{eq:graph-fiber-dimension}
 \dim V\bigl(I(W'),a,q+g\bigr)\leq k-2.
\end{equation}
The set is empty when $k=1$; the same convention applies if $V_W(a)$
is empty.
\end{lemma}

\begin{proof}
Write $A=\Qbar[W]$, $u=f/a^m$, and let $B$ be the finite normalization
of $A$.  The graph has coordinate ring $A[u]$, and $B[u]$ is finite
over it.  Suppose there were a point of the graph above a generic
point $\mathfrak p$ of a component of $V_W(a)$.  Lying over in
$B[u]$ and then contraction to $B$ give a prime $\mathfrak r$ above
$\mathfrak p$.  The local ring $B_{\mathfrak r}$ is a discrete
valuation ring.  The hypothesis says $u=v\pi^{-s}$ there, for a unit
$v$, a uniformizer $\pi$, and $s>0$.  Its map into the local ring of
the lifted graph point would satisfy $\pi^su=v$ with $\pi$ in the
maximal ideal and $v$ invertible, a contradiction.

The constructible image of $V_{W'}(a)$ therefore contains none of the
generic points of $V_W(a)$.  Each component of the latter has dimension
$k-1$, so the closure of this image has dimension at most $k-2$.
On imposing $q=-g$, projection from the graph to $W$ is a closed
immersion: in its coordinate ring the extra coordinate is already
specified by $g$.  Its image lies in the preceding constructible image,
which proves \eqref{eq:graph-fiber-dimension}.
\end{proof}

\subsection{Interpolation by successive powers of the parameter}

The extra equation $q+g=0$ in \Cref{lem:graph-fiber} is essential.
For example, on $W=\mathbb A^2$ with coordinates $(a,u)$, take $f=u$
and $m=1$.  The graph $W'=\{aq=u\}$ still has a one-dimensional
fiber over $a=0$: its entire $q$-line maps to $(a,u)=(0,0)$.
Imposing $q+g_0(u)=0$ leaves one point.  Thus it is the projected
boundary, and then the graph with this additional equation, that has
dimension $k-2$; the boundary of the graph alone need not have that
dimension.  \Cref{fig:graph-projection} displays this distinction.

\begin{figure}[H]
\centering
\begin{tikzpicture}[
  box/.style={draw,rounded corners,align=center,inner sep=6pt,
              text width=49mm,font=\small},
  >=Latex]
  \node[box] (graph) {$W'=\{u=aq\}$\\coordinates $(a,q)$};
  \node[box,right=21mm of graph] (base)
    {$W=\mathbb A^2$\\coordinates $(a,u)$};
  \node[box,below=14mm of graph] (fiber)
    {$W'\cap\{a=0\}$: the $q$-line\\dimension $1$};
  \node[box,below=14mm of base] (image)
    {image $(a,u)=(0,0)$\\dimension $0$};
  \draw[->] (graph) -- node[above,font=\small] {projection} (base);
  \draw[{Hooks[right]}->] (fiber) -- (graph);
  \draw[{Hooks[right]}->] (image) -- (base);
  \draw[->] (fiber) -- node[above,font=\small] {contracts} (image);
  \node[box,below=14mm of fiber] (point)
    {$a=0,\quad q=-g_0(0)$\\dimension $0$};
  \draw[{Hooks[right]}->] (point) --
    node[right,font=\small] {$q+g_0(u)=0$} (fiber);
\end{tikzpicture}

\caption{The graph $q=u/a$ over a two-dimensional affine base.
Horizontal arrows are projections and vertical arrows are inclusions.
The graph boundary is a line, but its image is a point.  The additional
equation $q+g_0(u)=0$ selects one point on that line, giving the
dimension $k-2=0$ used in the layer count.}
\label{fig:graph-projection}
\end{figure}

The next lemma gives a polynomial of low total degree with an arbitrarily
large prescribed polynomial ratio of order to degree.  Requiring
nonvanishing on the variety is part of the construction.

\begin{lemma}[Interpolation through successive layers]\label{lem:layer-interpolation}
Let $W'\subset\mathbb A^M$ be irreducible of dimension $k\geq1$,
with coordinates including $a,U,q$.  Let $\mathfrak b=(G_1,\ldots,G_t)$
have controlled generators and $V(\mathfrak b)=W'$.  Suppose
\[
 \dim V(\mathfrak b,a,F_0)\leq k-2,
 \qquad F_0=q+g_0(U),
\]
and put $F_n=F_0+\sum_{j=1}^{n-1}a^j g_j(U)$, where $m\ge1$ is an
integer and $\deg g_j\leq C(m+j+1)$.  Here $M,C$ are fixed and all other
complexities, including $m$, are polynomially bounded in $d$; the
coefficient heights through index $n$ are bounded by
$Q(n,d)\Psi(h)$.

For any prescribed positive integer $r$ polynomially bounded in $d$,
there are positive integers $L$ and $n=rL$, both polynomially bounded
in $d$, and a polynomial $P$ such that
\begin{equation}\label{eq:layer-membership}
 \deg P\leq L,\quad P\notin I(W'),\quad
 P=\sum_{\lambda=1}^t A_\lambda G_\lambda+A F_n+a^n B.
\end{equation}
The degrees of all the coefficients in this representation are
polynomial in $d$, and their joint affine height, including that of $P$,
is at most $Q'(d)\Psi(h)$.  The polynomial $Q'$ may depend on the
prescribed polynomial bound for $r$.
\end{lemma}

\begin{proof}
Fix a degree-compatible monomial order and a Gr\"obner basis of
$\mathfrak j=(\mathfrak b,a,F_0)$, together with representations of
its elements in these displayed generators.  Their degrees and heights
have the bounds recorded above.  For each reducible monomial, fix a
basis element to divide by.  Division in degree at most $T$ can then be
performed by a fixed decreasing sweep through all monomials of that
degree range, allowing zero coefficients.  Thus, for $\deg H\leq T$,
it gives linearly in $H$
\begin{equation}\label{eq:layer-division}
 H=\operatorname{rem}_{\mathfrak j}H
   +\sum_\lambda C_\lambda G_\lambda+aC+F_0D.
\end{equation}
There is a fixed bound $\Delta$, polynomial in the complexities at this
stage, such that every term on the right has degree at most $T+\Delta$.
Indeed, degree-compatible division produces quotient multipliers of
degree at most $T-\deg H_\alpha$ for a basis element $H_\alpha$.
Substituting its fixed generator representation adds its excess degree,
at most $\Delta$; it does not multiply $T$.

Work modulo $(\mathfrak b,F_n,a^n)$.  In
\eqref{eq:layer-division}, the $aC$ term advances one layer, while
\[
 F_0D\equiv-\sum_{j=1}^{n-1}a^jg_jD
\]
advances at least one layer.  Starting with a polynomial of degree at
most $L$, apply this division to each successive coefficient of
$1,a,\ldots,a^{n-1}$.  The coefficient being divided at layer $i$ has
degree at most $L+B_0(i+1)$, for a polynomially bounded $B_0$.
To verify this induction, a jump by $j\geq1$ adds at most
$\Delta+C(m+j+1)\leq B_0j$ to the degree.  The sum of all jumps to
layer $i$ is $i$, and there is no substitution at the same layer.
We retain every remainder and require it to be zero.

When $k=1$, the dimension hypothesis makes $\mathfrak j=(1)$, so there
are no remainder conditions.  If $\mathfrak j\ne(1)$, then $k\ge2$,
and the number of standard monomials of degree at most $T$ is at most
\[
 B_1(1+T)^{k-2},
\]
with $B_1$ polynomially bounded.  Here is an elementary uniform bound.
If the monomial initial ideal has generators of degree at most $D_1$,
in a standard monomial at most $\dim V(\mathfrak j)$ exponents can
be at least $D_1$.  Otherwise its variables of large exponent would
contain the support of a monomial generator, by the coordinate-prime
description of the radical of a monomial ideal.  Choose those variables
and bound every other exponent by $D_1-1$.  This gives the asserted
bound with a constant depending polynomially on $D_1$, since $M$ is
fixed.  If $\mathfrak j=(1)$ there are no standard monomials and no
conditions.  Thus at most
\begin{equation}\label{eq:layer-condition-count}
 B_1n(1+L+B_0n)^{k-2}
\end{equation}
linear conditions suffice when the ideal is proper.

Choose $k$ algebraically independent functions among the ambient
coordinate functions on $W'$.  Such a subset exists because these
functions generate its function field.  Let $V_L$ be the span of their
monomials of total degree at most $L$, regarded as actual polynomials
in the ambient coordinates.  Then
\[
 \dim V_L=\binom{L+k}{k},\qquad V_L\cap I(W')=0.
\]
This chooses independent polynomial representatives; no remainder map
is presumed to descend canonically to classes modulo $I(W')$.
One can equally extend $V_L$ to a linear complement of
$I(W')\cap\Qbar[a,U,q]_{\leq L}$.

Set $n=rL$.  For a proper $\mathfrak j$, the bound
\eqref{eq:layer-condition-count} is at most
\[
 B_2r(1+B_0r)^{k-2}L^{k-1}
\]
for a polynomially bounded $B_2$.  A sufficiently large polynomial
choice of $L$ makes it smaller than $\binom{L+k}{k}$.  In the unit-ideal
case the conclusion holds already for any nonzero source vector.
There is consequently a nonzero element $P\in V_L$ for which all
remainders vanish.  This proves both membership and $P\notin I(W')$.

For completeness, the construction also gives the asserted affine
height bounds.  Division in bounded degree is a linear operation on
a coefficient space of polynomial dimension.  During division the
leading monomial strictly decreases among the monomials of degree at
most $T$; thus its coefficient transformation is a product of
polynomially many matrices with controlled affine heights.  For
composable matrices with inner dimension $b$, the coefficient-tuple
convention gives
\[
 \ha{AB}\leq\ha{A}+\ha{B}+\log b.
\]
Indeed, locally every entry is a sum of $b$ products, and only infinite
places contribute the factor $b$.  The $n$ layer operations therefore have
polynomial affine heights.
Their source basis consists of monomials, so has height zero.  A kernel
vector obtained by minors and normalized to have a nonzero coordinate
equal to $1$ has affine height $Q(d,n,L)\Psi(h)$.

The tracked divisions give a membership representation of bounded
degree.  Alternatively, apply the arbitrary-membership bound to
$(\mathfrak b,F_n,a^n)$ and solve its coefficient equations with the
chosen $P$: the degrees and number of unknowns are polynomial in
$d,n,L$, and minors give the same affine height estimate.  This bounds
the scale of $P,A_\lambda,A,B$ together.  Since $r,L,n$ are polynomial
in $d$, \eqref{eq:layer-membership} follows with the stated bounds.
\end{proof}

\subsection{Product formula and descent}

\begin{proof}[Proof of \Cref{thm:interpolation}]
We descend through irreducible subvarieties $W\subseteq X$ containing
$z_*$.  At every step their degrees and degrees of chosen
set-theoretic equations are polynomial in $d$, and the affine heights
of those equations are at most $Q(d)\Psi(h)$.  Initially take a
component of $X$ containing the point.  The algebraic bounds above
preserve these conditions under each operation below.

If $a$ is constant on $W$, its constant value has height at most
$Q(d)\Psi(h)$, by projection to the $a$-line and the component height
bound.  On the other hand the product formula and
\eqref{eq:interpolation-smallness} give
\begin{equation}\label{eq:constant-parameter-height}
 \ha{a_*}=\sum_vw_v\log\max(1,\abs{a_*}_v^{-1})
 \geq S_{\rm small}\geq h/Q_1(d).
\end{equation}
Hence $h\leq Q(d)\Psi(h)$, which implies a polynomial bound in $d$:
for fixed $c$, $1+(\log h)^c\leq C_c\sqrt h$ for $h\geq2$.

Suppose now that $a$ is nonconstant and $k=\dim W\geq1$.
If $a$ is nonconstant but $V_W(a)$ is empty, this case can be settled
without an analytic equation.  For the controlled ideal
$\mathfrak b_W=(G_1,\ldots,G_t)$, the ideal $(\mathfrak b_W,a)$ is the
unit ideal.  The membership and coefficient bounds provide
\[
 1=\sum_i A_iG_i+aB
\]
with polynomial degrees in $d$ and joint affine height
$Q(d)\Psi(h)$.  At the target $1=a_*B(z_*)$.  Evaluating $B$ on the
selected embeddings costs only its coefficient height and the small
local sizes of $U_*$, since $\abs{a_*}_v<1$.  It follows that
\[
 S_{\rm small}=\sum_{v\in\mathcal V}w_v\log\abs{B(z_*)}_v
 \leq Q(d)\Psi(h).
\]
Again $h\leq Q(d)\Psi(h)$ gives the required polynomial bound.
We may therefore suppose that $V_W(a)$ is nonempty.

The total number of generic normalization branches above $V_W(a)$ is
at most $\deg W$.  Indeed take one general transverse linear slice
meeting every component $D_j$ of this fibre at general affine points.
The curve components have total degree at most $\deg W$.  Their
normalizations contain the transverse germs above all these points.
If a normalization divisor above $D_j$ has residue degree $f_i$ and
$e_i=\operatorname{ord}_i(a)$, it contributes
$\deg(D_j)f_ie_i$ to the zero divisor of $a$ on the normalized slice.
That divisor has degree equal to the polar degree of $a$, at most
$\deg W$.  In particular the number of branches is bounded.  This
counts residue degrees and ramification multiplicities, even though
the normalization itself is birational.

At each such branch choose the lowest nonzero order among the $E_\nu$.
To avoid cancellation of these leading terms in a combination, set
\[
 E=\sum_{\nu=1}^s t^{\nu-1}E_\nu
\]
for a positive integer $t$.  Each branch excludes at most $s-1$ values
of $t$, since its leading-coefficient polynomial is nonzero.  There is
therefore a permissible $t$ of polynomial size in $d$.  This combination preserves
the coefficient and tail bounds: at finite places its coefficients
are integral, and the extra triangle-inequality factor at infinite
places has total logarithm polynomial in $d$.  Absorb it in the $b_v$.
Write $E=\sum a^j e_j(U)$.

By \Cref{lem:zero-estimate}, an integer $m$ polynomially bounded in
$d$ can be chosen so that, for
\[
 f=\sum_{j=0}^{m-1}a^je_j(U),
\]
one has $\operatorname{ord}_{\mathfrak p}f
<m\operatorname{ord}_{\mathfrak p}a$ on every relevant branch.
Indeed the omitted terms have order at least
$m\operatorname{ord}_{\mathfrak p}a$, since all their coefficients
are regular on the affine variety.  Form the graph $W'$ of
$q=f/a^m$.  Its equations with controlled support are obtained by
saturating $(\mathfrak b_W,a^mq-f)$ by $a$; the closure of its support
is exactly the graph.  This takes place in one more variable and
has polynomial degree and affine height bounds.

The target lifts to $z_*'=(a_*,U_*,q_*)$.  Its global height is bounded
using the rational expression for $q_*$:
\begin{equation}\label{eq:graph-global-height}
 \ha{z_*'}\leq Q_5(d)h+Q_6(d)\Psi(h).
\end{equation}
Here $Q_5$ is fixed at this stage, before the interpolation degree and
order are chosen.  Its selected local size must instead be estimated
using the exact equation $E(z_*)=0$.  The tail bound gives
\begin{equation}\label{eq:graph-local-height}
 \abs{q_*}_v
 =\left|a_*^{-m}\sum_{j\geq m}a_*^je_j(U_*)\right|_v
 \leq\exp((m+1)b_v).
\end{equation}
In particular the sum of the positive local logarithmic sizes of
$(U_*,q_*)$ is at most $Q_7(d)\Psi(h)$.
The new equation and its finite truncations are
\[
 F=q+\sum_{j\geq0}a^je_{m+j}(U),\qquad
 F_n=q+\sum_{j=0}^{n-1}a^je_{m+j}(U).
\]
At the target,
\begin{equation}\label{eq:shifted-tail}
 \abs{F_n(z_*')}_v
 \leq\abs{a_*}_v^n\exp((m+n+1)b_v).
\end{equation}

By \Cref{lem:graph-fiber},
$V(I(W'),a,q+e_m)$ has dimension at most $k-2$.
For a controlled ideal $\mathfrak b$ defining $W'$ set-theoretically,
$V(\mathfrak b,a,q+e_m)$ is the same set and has the same dimension.
Apply \Cref{lem:layer-interpolation}, with the ratio $r=n/L$ to be
chosen shortly.  At the lifted target its identity becomes
\[
 P(z_*')=A(z_*')F_n(z_*')+a_*^nB(z_*').
\]
Equations \eqref{eq:polynomial-evaluation},
\eqref{eq:graph-local-height}, and \eqref{eq:shifted-tail}, together
with the coefficient and membership degree bounds, imply
\begin{equation}\label{eq:local-interpolation-value}
 \sum_{v\in\mathcal V}w_v\log\abs{P(z_*')}_v
 \leq -n S_{\rm small}+Q_8(d,n,L)\Psi(h).
\end{equation}
All positive contributions here are sums of affine coefficient heights,
positive local coordinate sizes, the $b_v$, and infinite-place
triangle-inequality terms.  There is no error proportional to the
number of finite places.
At the remaining places evaluate $P$ itself, rather than its membership
representation.  Since $\deg P\leq L$, \eqref{eq:graph-global-height}
gives
\begin{equation}\label{eq:outside-interpolation-value}
 \sum_{v\notin\mathcal V}w_v\log\abs{P(z_*')}_v
 \leq LQ_5(d)h+Q_9(d,n,L)\Psi(h).
\end{equation}

The choices can now be made in the required order.  The graph and
$m$ have already fixed $Q_5$.  Choose an integer
$r>2Q_1(d)(Q_5(d)+1)+1$, polynomially bounded in $d$.
Next choose the polynomially large $L$ supplied by
\Cref{lem:layer-interpolation}, and set $n=rL$.
If $P(z_*')\ne0$, its product formula and
\eqref{eq:local-interpolation-value}--\eqref{eq:outside-interpolation-value}
give
\[
 0\leq
 \left(-\frac n{Q_1(d)}+LQ_5(d)\right)h
 +Q_{10}(d,n,L)\Psi(h).
\]
The coefficient of $h$ is at most $-L$.  As $n,L$ are now fixed
polynomials in $d$, this inequality
again implies $h\leq Q(d)\Psi(h)$ and hence a polynomial height
bound.  Thus, beyond a polynomial threshold chosen \emph{after}
$m,r,L,n$, the value $P(z_*')$ must be zero.

Since $P\notin I(W')$, the intersection $W'\cap V(P)$ has dimension
at most $k-1$.  Project it to the original $(a,U)$-space, take its
closure, and choose an irreducible component containing $z_*$.  Such
a component exists because $a_*\ne0$ and the graph projection is an
isomorphism over that open set.  Projection cannot increase dimension.
The degree and arithmetic bounds above preserve the induction
hypotheses.  Discard the graph coordinate $q$; at the next step construct
a new graph from the original common equations on this smaller
variety.  Consequently the ambient number of variables stays fixed.

The dimension drops at every continuing step, so there are at most
$\dim X$ such steps.  A zero-dimensional final variety has constant
$a$ and is covered by \eqref{eq:constant-parameter-height}.  Taking the
maximum of the finitely many polynomial thresholds proves the theorem.
All component and coefficient field extensions made during the proof
extend the selected local data by all embeddings, as specified above;
they do not alter any smallness sum or introduce a field-degree factor.
\end{proof}
